% =====================================================================
%  Conclusion générale
% =====================================================================
\chapter*{Conclusion générale}
\addcontentsline{toc}{chapter}{Conclusion générale}
\markboth{Conclusion générale}{Conclusion générale}

Réalisé au sein de France Téléphone, filiale du groupe Winvest Capital, ce
projet de fin d'études avait pour ambition de concevoir et de développer,
de bout en bout, une plateforme SaaS multi-entreprises de création de
modèles de communication assistée par intelligence artificielle. Au terme
de ce travail, l'objectif fixé a été atteint : \textbf{Winaity Template
Builder} est une plateforme fonctionnelle et déployée en production.

\paragraph{Récapitulation de la démarche.}
La conduite du projet a suivi la méthodologie agile Scrum, structurée en
quatre sprints. Après une phase d'analyse des besoins et de conception —
identification des acteurs, expression des besoins, modélisation UML et
définition de l'architecture microservices —, nous avons successivement mis
en \oe{}uvre l'authentification et le modèle multi-entreprises (sprint~1),
l'éditeur de modèles multi-canal (sprint~2), l'assistant d'intelligence
artificielle agentique (sprint~3), puis les abonnements, les intégrations et
la mise en production (sprint~4).

\paragraph{Résultats obtenus.}
La plateforme répond à la problématique posée en introduction. Elle permet à
une équipe métier, sans compétence technique, de produire rapidement des
supports sur cinq canaux (e-mail, facture, contrat, SMS, RCS) au moyen d'un
éditeur visuel. Son assistant d'IA génère et modifie de véritables blocs
éditables et sait reconstruire une campagne à partir d'une simple affiche.
L'isolation multi-entreprises garantit la confidentialité des données de
chaque société, tandis que l'auto-hébergement du modèle de langage assure la
souveraineté des données — deux exigences qu'aucune des solutions du marché
étudiées ne réunissait. La plateforme intègre enfin la gestion des rôles, un
système d'abonnements Stripe conforme sur le plan fiscal, une couche
d'intégration, et un déploiement sécurisé.

\paragraph{Difficultés rencontrées.}
La principale difficulté a résidé dans la fiabilisation de l'assistant
d'intelligence artificielle. Un modèle de langage de taille modérée, servi
localement, se révèle peu fiable lorsqu'on le laisse enchaîner librement de
nombreux appels d'outils : duplications, omissions ou actions inopérantes.
Cette difficulté a été surmontée par une architecture hybride associant un
planificateur déterministe contraint par schéma, une voie agentique de
repli, un routeur d'intention et des mécanismes de fiabilité (garde de
contraste, confirmations honnêtes, filtrage des sorties du modèle).
L'exécution multimodale (lecture d'affiche) a également imposé d'isoler
l'étape de vision et de la fiabiliser par décodage contraint. Sur le plan de
l'ingénierie, la coordination d'une architecture microservices et sa mise en
production ont nécessité une attention particulière à la configuration, aux
secrets et à la reproductibilité des environnements.

\paragraph{Apports du stage.}
Ce projet a été l'occasion d'acquérir et d'approfondir de nombreuses
compétences. Sur le plan technique : la maîtrise d'une architecture
microservices moderne (NestJS, gRPC, CQRS, TypeORM, PostgreSQL, MinIO), du
développement frontend avec Next.js et React, de l'intégration d'un système
de paiement, et surtout de l'ingénierie appliquée aux grands modèles de
langage — génération agentique par outils, planification contrainte par
schéma, vision par ordinateur et protocole MCP, sur une infrastructure
d'inférence auto-hébergée. Sur le plan professionnel et personnel, ce stage
a permis d'évoluer au sein d'une équipe, de dialoguer avec des interlocuteurs
métier, de conduire un projet de son analyse à sa mise en production, et de
prendre des décisions d'architecture argumentées.

\paragraph{Perspectives.}
Plusieurs pistes d'évolution ont été identifiées. Sur le plan collaboratif,
l'introduction de fonctionnalités temps réel (curseurs et synchronisation des
brouillons) enrichirait le travail à plusieurs. Sur le plan de l'IA, le
recours à un modèle plus puissant pour la seule étape de planification
améliorerait la fidélité des mises en page complexes, et l'envoi effectif des
messages RCS via un opérateur compléterait le canal correspondant. Sur le
plan de l'écosystème, l'approfondissement du flux d'intégration permettrait
de connecter la plateforme à des outils tiers tels que des CRM, et le site
marketing pourrait être enrichi (contenu, référencement, démonstrations). Ces
perspectives confirment le potentiel de la plateforme et ouvrent la voie à
son développement continu.
